% LSPV-5 D3CAL1 — 2026-10-02
\section*{D3CAL1: Three-variable clause-count saturation control}
\textbf{STATUS: PROVED/EXACT FINITE CONTROL.}

With three variables and simple exact-3 clauses using three distinct variables, there are exactly eight possible signed clauses. Each clause falsifies one distinct assignment in \(\{0,1\}^3\).

Hence an exact-m formula with \(m<8\) distinct clauses leaves \(8-m>0\) satisfying assignments. Every such formula is SAT. For \(m=8\), the only formula contains all eight clauses and is UNSAT. For \(m=9\), a simple distinct-clause formula does not exist; allowing duplicates gives no new semantic constraint because \(C\wedge C=C\).

Therefore, under the promise that the input belongs to the exact \(D(3,m)\) family, the minimum SAT/UNSAT decision object is a constant SAT terminal for \(m<8\) and a constant UNSAT terminal for \(m=8\). This corrects an earlier bookkeeping count: the 8-node ordered membership DAG for \(D(3,8)\) represented the process of selecting the unique eight-clause set, not intrinsic SAT-solving complexity.

This is a useful hostile control. Formula-enumeration machinery must not be counted as solver complexity when the promise itself determines the answer.

The first nontrivial fixed-variable clause-count experiment is therefore \(D(4,m)\), where SAT and UNSAT formulas coexist beginning at the cube-edge-cover threshold.

P versus NP remains OPEN.
% END_OF_NOTE::D3CAL1
